% OWNH whitepaper. Build: npm run whitepaper (writes docs/ownh-whitepaper.pdf).
% Describes the method as implemented in src/; keep it in step with the code.
% Deadpan pitch voice, like docs/index.html. No figures from runs.
\pdftrailerid{} % no build-path-dependent PDF ID: same source, same bytes
\documentclass[11pt,a4paper]{article}
\usepackage[T1]{fontenc}
\usepackage[utf8]{inputenc}
\usepackage{charter}
\usepackage[scaled=0.92]{helvet}
\usepackage{lmodern}
\renewcommand{\rmdefault}{bch}
\usepackage[margin=2.6cm]{geometry}
\usepackage{microtype}
\usepackage{amsmath,amssymb}
\usepackage{booktabs,tabularx}
\usepackage[hyphens]{url}
\usepackage{xcolor}
\usepackage{fancyhdr}
\usepackage[font=small,labelfont=bf]{caption}
\definecolor{brand}{HTML}{1F4FD1}
\usepackage[colorlinks,linkcolor=brand,urlcolor=brand,citecolor=brand,
  pdftitle={Content-Addressable Accountability: A Line-Hash Approach to Code Ownership},
  pdfauthor={OWNH},pdfsubject={The OWNH method},pdfcreator={OWNH}]{hyperref}

\pagestyle{fancy}
\fancyhf{}
\fancyhead[L]{\sffamily\footnotesize\textcolor{brand}{\textbf{OWNH}}\ Whitepaper}
\fancyhead[R]{\sffamily\footnotesize Objective Weighted Normalized Heuristics}
\fancyfoot[C]{\sffamily\footnotesize\thepage}
\renewcommand{\headrulewidth}{0.4pt}
\setlength{\headheight}{14pt}
\setlength{\parskip}{0.45em}
\setlength{\parindent}{0pt}

\makeatletter
\renewcommand\section{\@startsection{section}{1}{\z@}{-2.2ex plus -0.5ex}{1ex}{\sffamily\large\bfseries\color{brand}}}
\renewcommand\subsection{\@startsection{subsection}{2}{\z@}{-1.6ex plus -0.4ex}{0.6ex}{\sffamily\normalsize\bfseries}}
\makeatother

\newcommand{\ownh}{OWNH}
\newcommand{\key}[1]{\texttt{#1}}
\newenvironment{definition}[1]{\par\smallskip\noindent\textbf{Definition (#1).}\ }{\par\smallskip}
\newenvironment{property}[1]{\par\smallskip\noindent\textbf{Property (#1).}\ }{\par\smallskip}

\begin{document}

\thispagestyle{empty}
\begin{center}
  {\sffamily\small\color{brand}\textbf{OWNH WHITEPAPER}}\par\vspace{1.2em}
  {\LARGE\bfseries Content-Addressable Accountability:\\[0.2em] A Line-Hash Approach to Code Ownership\par}
  \vspace{1em}
  {\large The \ownh{} Framework\par}
  \vspace{0.3em}
  {\small\sffamily Objective Weighted Normalized Heuristics\par}
  \vspace{0.3em}
  {\small\url{https://ownh.org}\par}
\end{center}
\vspace{1em}

\begin{abstract}
\noindent Engineering organizations need to know who owns their code. Existing
answers depend on activity (commit counts), recency (\texttt{git blame}), or
declarations (\texttt{CODEOWNERS} files), and each can be changed by the
behavior it is meant to measure. This paper presents \ownh{}, a method that
assigns every line of code to exactly one owner using only the line's content
hash and the version history of the repositories that contain it. The owner of
a line is the first contributor, across all repositories under analysis, to
commit a line with that exact content. The assignment is deterministic,
independent of the order in which repositories are analyzed, and permanent for
a fixed scope of analyzed histories.
The definitions, the attribution procedure, the benchmark against
established methods, and the derived governance metrics: ownership leverage,
code demolition, code survival, and quarterly projections suitable for
objectives and key results, are described.
\end{abstract}

\section{Introduction}

Ownership measures have a respectable research history. In Windows Vista and
Windows 7, ownership measures correlated with pre-release faults and
post-release failures~\cite{bird}. Research comparing ownership definitions
finds that line-based ownership suits accountability and attribution, that
commit-based ownership is the better choice for bug fixing and quality
planning, and that commit-based and line-based methods disagree about who the
major developers are~\cite{lines}. The question is therefore not whether
ownership matters, but which definition an organization can defend. The method
presented here addresses accountability.

Each conventional definition carries a human element:

\begin{itemize}\setlength{\itemsep}{0.2em}
  \item \textbf{Commit counts} reward volume. Ten small commits outrank one good one.
  \item \textbf{\texttt{git blame}} reports the revision and author that last
    modified each line~\cite{blame}. Recency is not authorship.
  \item \textbf{\texttt{CODEOWNERS}} files are declarations. They are maintained
    by people and are out of date by the time they merge.
\end{itemize}

In the method presented here, the human element is removed. The history is
read once, and a single answer is produced for every line, which any party can
regenerate from the same repositories.

\section{Definitions}

\begin{definition}{Line}
A line is the exact sequence of bytes between two newline characters. No
normalization is applied. Leading and trailing whitespace, indentation, blank
lines, and a trailing carriage return are part of the line. A line ending in
\verb|\r\n| and the same line ending in \verb|\n| are different lines.
\end{definition}

\begin{definition}{Binary file}
A file that git classifies as binary, using its attributes and content check,
is treated as a single line whose content is the entire file. Git's
classification is used as is; no other heuristic is applied.
\end{definition}

\begin{definition}{Line hash}
The identity of a line is $h = \mathrm{SHA\text{-}256}(\ell)$, the hash of its
bytes. Two lines are the same line if and only if their hashes are equal,
regardless of file, path, language, or repository.
\end{definition}

\begin{table}[h]
\centering
\small
\begin{tabular}{@{}lll@{}}
\toprule
Line (shown between brackets) & SHA-256 prefix & Remark \\
\midrule
\texttt{[]}             & \texttt{e3b0c44298fc1c14} & a blank line \\
\texttt{[\char`\}]}     & \texttt{d10b36aa74a59bcf} & closing brace \\
\texttt{[\ \ \char`\}]} & \texttt{737db166c79ae98e} & closing brace, one level in \\
\texttt{[\ \ \ \ \char`\}]} & \texttt{28d86778615f6af4} & closing brace, two levels in \\
\texttt{[\textbackslash r]} & \texttt{9d1e0e2d9459d065} & a blank line, CRLF \\
\bottomrule
\end{tabular}
\caption{Line hashes. Each row is a distinct line with its own owner.}
\end{table}

\begin{definition}{Introduction}
A commit $c$ introduces hash $h$ if the diff of $c$ against its parent adds a
line with hash $h$. Each introduction carries the key
\[
  k(c) = \bigl(t_{\mathrm{author}}(c),\ \mathrm{repository}(c),\ \mathrm{position}(c)\bigr),
\]
where $t_{\mathrm{author}}$ is the author time, $\mathrm{repository}$ is the
name of the repository, unique within an analysis, and $\mathrm{position}$ is
the commit's index in the canonical traversal of that repository's history:
git's topological order, oldest commit first. A history admits several
topological orders; git computes this one deterministically for a given
history.
\end{definition}

\begin{definition}{Owner}
The owner of hash $h$ is the author of the introduction with the smallest key:
\[
  \mathrm{owner}(h) = \mathrm{author}\Bigl(\operatorname*{arg\,min}_{c \,\in\, I(h)} k(c)\Bigr),
\]
where $I(h)$ is the set of commits, in any repository under analysis, that
introduce $h$. Every occurrence of $h$, in every file of every repository, has
this owner.
\end{definition}

\begin{definition}{Identity}
Contributors are identified by email address, compared without regard to case,
after applying each repository's \texttt{.mailmap}. The display name of an
identity is the name it has used most often.
\end{definition}

\section{Attribution}

\subsection{Introductions and removals}
The complete history of every repository is considered. Each line a commit adds
relative to its parent is an introduction of that line's hash, and the owner of
each hash follows from the smallest key among its introductions. Each line a
commit removes is a removal, recorded by hash, quarter, and the identity of the
commit's author, for the analyses in Section~\ref{sec:metrics}.

\subsection{Merge commits}
A merge commit is compared with the result of merging its parents
automatically. The lines it adds relative to that result are the lines its
author resolved by hand, and they are introductions like any other. Merge
commits are not counted in commit totals, and lines they remove are not
counted as removals. As evidence of activity, a merge counts like any other
commit.

\subsection{Lines at head}
Every line of every file at the head of each repository is attributed to the
owner of its hash. The total is the number of \emph{lines under management}.
Each line counts once per occurrence: a hash that appears in many places
contributes all of those places to its owner.

\section{Properties}

\begin{property}{Uniqueness}
Every line has exactly one owner. Ties are impossible: two introductions with
equal author time are separated by repository name, and within a repository by
position in the canonical traversal.
\end{property}

\begin{property}{Determinism}
The same repositories at the same commits produce the same owners and the same
counts.
\end{property}

\begin{property}{Order independence}
Ownership is a minimum over a key that does not depend on the order of
analysis. The result does not depend on the order in which repositories are
analyzed.
\end{property}

\begin{property}{Permanence}
An owner keeps ownership regardless of later activity, inactivity, or
employment status. Ownership is permanent for a fixed scope: for a given set of
analyzed histories, ownership of a line changes only when the line changes.
Extending the scope with a repository that introduces the same line earlier,
or rewriting an analyzed history, transfers ownership to the earlier
introduction.
\end{property}

\begin{property}{Global scope}
Ownership is shared across repositories. A line first written in one
repository is owned by its first author in every repository where it occurs.
\end{property}

\subsection{Operational consequences}
The following follow directly from the definitions and are reported as
designed:

\begin{itemize}\setlength{\itemsep}{0.2em}
  \item \textbf{Formatting.} A formatting commit changes the hash of every line
    it reformats. Lines whose new form has not been written before are
    introduced by the author of the formatting commit, who becomes their owner.
  \item \textbf{Renames.} Renaming an identifier changes every line that
    contains it, with the same effect.
  \item \textbf{Reverts.} A revert restores earlier lines. Their hashes already
    have owners, so ownership returns to the original authors.
  \item \textbf{Indentation.} Each indentation level of a closing brace is a
    separate line with its own owner.
  \item \textbf{Copying.} A copied line keeps the owner it already has, wherever
    it is pasted.
  \item \textbf{Blank lines.} The blank line is one line. It has one owner.
\end{itemize}

\section{Benchmarking against established methods}

For each repository, and for all repositories combined, the top owner is
determined under three definitions:

\begin{enumerate}\setlength{\itemsep}{0.2em}
  \item \textbf{Commits}: the identity with the most non-merge commits.
  \item \textbf{Blame}: the identity with the most lines at head according to
    \texttt{git blame}.
  \item \textbf{Line hash}: the identity with the most lines at head by
    first introduction.
\end{enumerate}

For each repository, agreement is recorded when the three definitions name the
same owner.

Where a repository declares ownership in a \texttt{CODEOWNERS} file, each line
at head is assigned to the last rule that matches its file, and for each rule
the declared owners are set beside the top owner of the same lines by line
hash and by blame. Declared owners are teams or accounts rather than
contributors, so no agreement is recorded; the comparison is presented as
found.

\section{Governance metrics}
\label{sec:metrics}

\subsection{Ownership leverage}
Leverage is the number of lines an identity owns per line it wrote, where lines
written are the lines added in that identity's own commits:
\[
  \mathrm{leverage}(o) = \frac{\text{lines owned at head}}{\text{lines added in own commits}}.
\]
A leverage above one
means an identity's lines have been reproduced, by others or by itself, more
often than they were written.

\subsection{Code demolition}
A removal is attributed to the author of the removing commit and counted
against the owner of the removed line. From the recorded removals are derived:
the lines each identity removed that were owned by someone else, the lines each
owner lost to others, and the lines added again after they had first been
removed. A restored line returns to its original
owner.

\subsection{Code survival}
Each added line is followed from the quarter it was added. Copies of the same
hash cannot be told apart, so within each repository removals are paired with
the oldest surviving copies first. A paired line dies at its age in quarters; a line still present
is censored at its current age. Survival is the Kaplan--Meier estimate over
these ages, overall and per owner, in the tradition of cohort analyses of
surviving code~\cite{theseus}. The \emph{code half-life} is the first age at
which the estimate falls to one half or below. When the estimate never falls
that far, the half-life is extrapolated with an exponential through the last
point and labeled as projected.

\subsection{Outlook}
The ownership series is the cumulative net of recorded additions and removals
per owner and quarter, summed over all branches of the history. It is not a
sequence of snapshots: removals inside merge commits are not counted, so lines
added on a branch whose changes a merge later discarded remain in the series.
The difference from the lines at head is stated with each outlook. ``Now'' is
the quarter of the newest commit. An identity is
\emph{inactive} in a quarter if it has no commit, merges included, in that
quarter or the three before it. Three series are projected:

\begin{itemize}\setlength{\itemsep}{0.2em}
  \item the share of lines owned by the principal owner;
  \item the share of lines owned by inactive identities;
  \item the number of blank lines.
\end{itemize}

Each projection is a least-squares line fitted over the last 12 quarters, with
at least three quarters required, continued from the latest value and reported
with its coefficient of determination $R^2$. For a share, the projected result
is the quarter in which it reaches one half or more, or the quarter since which
it has held that. For blank lines, it is the quarter in which the next
milestone is reached, where milestones are the round numbers
$1, 2, 5, 10, 20, 50, \dots$ A series whose fitted trend is flat or falling has
no projected quarter, and a result beyond 40 quarters is stated as not expected
within that horizon.

\section{Related work}

Studies built on World of Code identify the repository that first committed
each version of a file as its origin~\cite{woc}. The Software Heritage
provenance service answers where a file comes from with the oldest revision
that contains it~\cite{swh}. Here the same first-introducer principle is
applied at the granularity of a single line. Truck-factor research derives key
developers from version history~\cite{truck}. \texttt{git blame} offers
options to follow moved and copied lines and to ignore formatting
revisions~\cite{blame}; in this method every revision is credited as written.

\section{Conclusion}

Under \ownh{}, the question of code ownership is settled by a definition that
requires no judgment: the first contributor to write a line owns it,
everywhere and permanently. Every line has an owner. Every owner can be named. Someone has to
own it.

\begin{thebibliography}{9}\small
\bibitem{bird} C.~Bird, N.~Nagappan, B.~Murphy, H.~Gall, P.~Devanbu.
  \emph{Don't Touch My Code! Examining the Effects of Ownership on Software Quality.}
  ESEC/FSE 2011.
  \url{https://www.microsoft.com/en-us/research/publication/dont-touch-my-code-examining-the-effects-of-ownership-on-software-quality/}
\bibitem{lines} \emph{Code Ownership: The Principles, Differences, and Their
  Associations with Software Quality.} arXiv:2408.12807.
  \url{https://arxiv.org/abs/2408.12807}
\bibitem{blame} \emph{git-blame documentation.} \url{https://git-scm.com/docs/git-blame}
\bibitem{woc} M.~Jahanshahi, D.~Reid, A.~Mockus.
  \emph{Beyond Dependencies: Copy-Based Reuse.}
  arXiv:2409.04830. \url{https://arxiv.org/abs/2409.04830}
\bibitem{swh} Software Heritage. \emph{Provenance service documentation.}
  \url{https://docs.softwareheritage.org/devel/swh-provenance/}
\bibitem{truck} G.~Avelino, M.~T.~Valente, et al.
  \emph{What is the Truck Factor of popular GitHub applications?}
  \url{https://peerj.com/preprints/1233}
\bibitem{theseus} \emph{git-of-theseus.} \url{https://github.com/erikbern/git-of-theseus}
\end{thebibliography}

\end{document}
